\section{Light --- Syntax}
%%% generated by the blueprint pipeline from TCSlib.Programs.Light.Syntax
\subsection*{Light --- Syntax: Introduction}

Declarations of the Lean module \texttt{TCSlib.Programs.Light.Syntax}.

\subsection*{Light --- Syntax: Results}

\begin{definition}[Arithmetic operations]
\lean{Light.Op}
\label{Light.Op}
\leanok
The arithmetic operations are addition, subtraction, and multiplication.
\end{definition}

\begin{definition}[Evaluation of integer operations]
\lean{Light.Op.eval}
\label{Light.Op.eval}
\leanok
\uses{Light.Op}
For integers $a$ and $b$, define the evaluations of addition, subtraction, and
multiplication to be $a+b$, $a-b$, and $ab$, respectively.
\end{definition}

\begin{definition}[Arithmetic expression syntax]
\lean{Light.Expr}
\label{Light.Expr}
\leanok
\uses{Light.Op}
An arithmetic expression is a natural-number constant, a variable indexed by a natural
number, a binary application of addition, subtraction, or multiplication to two
expressions, or a load from an address expression.
\end{definition}

\begin{definition}[Conditions]
\lean{Light.Cond}
\label{Light.Cond}
\leanok
\uses{Light.Expr}
A condition is either a strict comparison between two expressions or an equality between
two expressions.
\end{definition}

\begin{definition}[Imperative statement syntax]
\lean{Light.Stmt}
\label{Light.Stmt}
\leanok
\uses{Light.Cond, Light.Expr}
An imperative statement is built from a no-operation statement, assignment of an
arithmetic expression to a natural-number variable, storage of one expression at the
address given by another, sequential composition, a conditional with two branches, a
while loop, or a procedure call with a natural-number identifier, a list of expression
arguments, and a natural-number destination variable. Expressions are formed from
natural-number constants and variable indices, binary addition, subtraction, or
multiplication, and memory loads; conditions compare two expressions using either strict
inequality or equality.
\end{definition}

\begin{definition}[Programs as statement lists]
\lean{Light.Program}
\label{Light.Program}
\leanok
\uses{Light.Stmt}
A program is a finite list of imperative statements.
\end{definition}

\begin{definition}[Program state]
\lean{Light.State}
\label{Light.State}
\leanok
A state assigns an integer value to each natural-number variable index and an integer
value to each natural-number memory address.
\end{definition}

\begin{definition}[Resource limits]
\lean{Light.Limits}
\label{Light.Limits}
\leanok
A resource limit consists of an integer word bound, a natural-number space bound, and a
natural-number depth bound.
\end{definition}

\begin{definition}[Expression evaluation]
\lean{Light.Expr.val}
\label{Light.Expr.val}
\leanok
\uses{Light.Expr, Light.Op.eval, Light.State}
Given a state, evaluate an arithmetic expression as an integer: a constant has its
natural-number value, a variable has the value assigned to its index, and an operation
applies addition, subtraction, or multiplication to the values of its two
subexpressions. A load evaluates its address expression, converts that integer to a
natural number, and returns the memory value at that address.
\end{definition}

\begin{definition}[Expression cost]
\lean{Light.Expr.cost}
\label{Light.Expr.cost}
\leanok
\uses{Light.Expr}
The cost of an expression counts its syntax nodes: each constant or variable has cost
one; a binary operation has cost one plus the costs of its two subexpressions; and a
load has cost one plus the cost of its address expression.
\end{definition}

\begin{definition}[Valid memory address]
\lean{Light.Limits.Addr}
\label{Light.Limits.Addr}
\leanok
\uses{Light.Limits}
For resource limits with space bound $S$, an integer $a$ is a valid address exactly when
$0 \le a < S$.
\end{definition}

\begin{definition}[Expression safety]
\lean{Light.Expr.Safe}
\label{Light.Expr.Safe}
\leanok
\uses{Light.Expr, Light.Expr.val, Light.Limits, Light.Limits.Addr, Light.Op.eval, Light.State}
Given resource limits and a state, an expression is safe if it satisfies the following
recursive conditions. A constant is safe when its value is at most the integer word
bound, and every variable is safe. A binary operation is safe when both subexpressions
are safe and the absolute value of their result is at most the word bound; a load is
safe when its address expression is safe and evaluates to a valid address.
\end{definition}

\begin{definition}[Condition satisfaction]
\lean{Light.Cond.Holds}
\label{Light.Cond.Holds}
\leanok
\uses{Light.Cond, Light.Expr.val, Light.State}
A condition holds in a state when its two evaluated expressions satisfy the specified
comparison: the first is strictly less than the second for a less-than condition, or the
two values are equal for an equality condition.
\end{definition}

\begin{definition}[Condition cost]
\lean{Light.Cond.cost}
\label{Light.Cond.cost}
\leanok
\uses{Light.Cond, Light.Expr.cost}
The cost of a condition is the sum of the costs of its two expressions plus one, whether
the condition tests strict inequality or equality.
\end{definition}

\begin{definition}[Condition safety]
\lean{Light.Cond.Safe}
\label{Light.Cond.Safe}
\leanok
\uses{Light.Cond, Light.Expr.Safe, Light.Limits, Light.State}
A condition is safe under given limits and in a given state if both of its expressions
are safe under those limits and in that state, whether it tests strict inequality or
equality.
\end{definition}

\begin{definition}[Default-zero list frame]
\lean{Light.frame}
\label{Light.frame}
\leanok
For a list of integers, the frame is the function on natural-number indices that returns
the list entry at that index, or zero when the index is out of range.
\end{definition}

\begin{definition}[Default-zero memory function]
\lean{Light.memOf}
\label{Light.memOf}
\leanok
For a list of integers, the memory function is the function on natural-number addresses
that returns the list entry at that address, or zero when the address is out of range.
\end{definition}

\begin{definition}[Bounded-depth execution relation]
\lean{Light.Exec}
\label{Light.Exec}
\leanok
\uses{Light.Cond.Holds, Light.Cond.Safe, Light.Cond.cost, Light.Expr.Safe, Light.Expr.cost, Light.Expr.val, Light.Limits, Light.Limits.Addr, Light.Program, Light.State, Light.Stmt, Light.frame}
For fixed limits and a program (a list of statements), the execution relation records
that a statement, started at depth $d$ in an initial state, terminates in a final state
with cost $k$. A skip leaves the state unchanged at cost zero; assignment and store
require safe expressions (and a valid address for a store) and update the corresponding
location or memory entry, at expression cost plus one. Sequential composition adds the
two execution costs. A conditional requires a safe condition and executes the branch
selected by whether the condition holds, at condition cost plus one plus the branch
cost. A loop requires a safe condition and either exits immediately when the condition
does not hold, at condition cost plus one, or executes its body and then the loop again
when the condition holds, adding both execution costs to the condition cost plus one. A
procedure call requires every argument expression to be safe, the program entry at index
$p$ to be the called body, and $d < 	ext{depth limit}$; it executes that body at depth
$d+1$ in a state whose locations are initialized from the evaluated arguments (with
missing entries zero) and whose memory is unchanged, then updates location $x$ with the
body's final location zero and adopts its final memory. Its cost is the sum of
argument-expression costs plus two plus the body's execution cost.
\end{definition}

\begin{definition}[Bounded state]
\lean{Light.State.Bounded}
\label{Light.State.Bounded}
\leanok
\uses{Light.Limits, Light.State}
A state is bounded under given limits if every location value and every memory value has
absolute value at most the word limit.
\end{definition}
